\section*{Hoofdstuk \ref{ch:g12:comb} --- Combinatoriek en tellen}

\begin{solution}{exo:g12:comb:1}
Productprincipe:
$26^2 \times 10^3 \times 26^2 = 26^4 \times 10^3 = 456\,976\,000$.

Zonder herhaald teken moeten de vier letters verschillend zijn
($26 \times 25 \times 24 \times 23$ manieren, door de letterplaatsen in
volgorde in te vullen) en de drie cijfers ook ($10 \times 9 \times 8$):
\[
26 \times 25 \times 24 \times 23 \times 10 \times 9 \times 8
= 358\,800 \times 720 = 258\,336\,000 .
\]
\end{solution}

\begin{solution}{exo:g12:comb:2}
$\dbinom83 = \dfrac{8 \times 7 \times 6}{3!} = 56$;
$\dbinom{10}{8} = \dbinom{10}{2} = \dfrac{10 \times 9}{2} = 45$;
\[
\frac{\binom n2}{\binom{n+1}2}
= \frac{n(n-1)/2}{(n+1)n/2} = \frac{n-1}{n+1}.
\]
\end{solution}

\begin{solution}{exo:g12:comb:3}
Kies het comité: $\binom{30}{4}$ manieren. Kies daarna de voorzitter en de
penningmeester uit de 4, in volgorde: $4 \times 3 = 12$ manieren. In totaal
\[
\binom{30}{4} \times 12 = 27\,405 \times 12 = 328\,860 .
\]
\end{solution}

\begin{solution}{exo:g12:comb:4}
\[
(x+2)^5 = x^5 + 10x^4 + 40x^3 + 80x^2 + 80x + 32 ,
\]
\[
(1-x)^6 = 1 - 6x + 15x^2 - 20x^3 + 15x^4 - 6x^5 + x^6 .
\]
In $(2x+3)^7$ is de term in $x^3$ gelijk aan
$\binom{7}{3}(2x)^3\,3^4 = 35 \times 8 \times 81\, x^3$: de coëfficiënt is
$22\,680$.
\end{solution}

\begin{solution}{exo:g12:comb:5}
\emph{1.} $\dbinom{52}{5} = 2\,598\,960$.

\emph{2.} Precies één aas: kies hem ($4$ manieren) en vul aan met $4$
niet-azen: $4 \times \binom{48}{4} = 4 \times 194\,580 = 778\,320$.
Minstens één aas: tel via het \omterm{def:g10:proba:operations}{complement},
$\binom{52}{5} - \binom{48}{5} = 2\,598\,960 - 1\,712\,304 = 886\,656$.

\emph{3.} Kies de waarde van het drietal ($13$), zijn kleuren
($\binom43 = 4$), de waarde van het paar ($12$ resterende), zijn kleuren
($\binom42 = 6$):
$13 \times 4 \times 12 \times 6 = 3744$.
\end{solution}

\begin{solution}{exo:g12:comb:6}
GETAL heeft 5 verschillende letters: $5! = 120$ anagrammen.

BANANA heeft 6 letters: drie A's, twee N's, één B. Kies de plaatsen van de
A's ($\binom63$), daarna die van de N's onder de rest ($\binom32$); de B
neemt de laatste plaats:
\[
\binom{6}{3}\binom{3}{2} = 20 \times 3 = 60 .
\]
(Even goed: $\frac{6!}{3!\,2!\,1!} = 60$.)
\end{solution}

\begin{solution}{exo:g12:comb:7}
\emph{Algebraïsch:}
\[
k\binom nk = \frac{k\,n!}{k!(n-k)!} = \frac{n!}{(k-1)!\,(n-k)!}
= n\,\frac{(n-1)!}{(k-1)!\bigl((n-1)-(k-1)\bigr)!} = n\binom{n-1}{k-1}.
\]
\emph{Met dubbel tellen:} tel de paren (comité van $k$ personen, voorzitter
erin). Ofwel kies je het comité ($\binom nk$) en daarna zijn voorzitter
($k$): $k\binom nk$ paren. Ofwel kies je eerst de voorzitter ($n$
mogelijkheden) en daarna de $k-1$ andere leden uit de $n-1$ overige:
$n\binom{n-1}{k-1}$ paren.
\end{solution}

\begin{solution}{exo:g12:comb:8}
Een pad bestaat uit precies $m + n$ stappen, waarvan er $m$ naar het oosten
en $n$ naar het noorden gaan; het ligt volledig vast door de verzameling
tijdstippen (uit de $m+n$) waarop je naar het oosten stapt. Er zijn
$\binom{m+n}{m}$ zulke keuzen.
\end{solution}

\begin{solution}{exo:g12:comb:9}
Splits een verzameling van $m + n$ personen in een groep $A$ van $m$ en een
groep $B$ van $n$. Een deelverzameling met $k$ elementen bevat een zeker
aantal $j$ leden van $A$ ($0 \leq j \leq k$) en $k - j$ leden van $B$; voor
vaste $j$ zijn er $\binom mj \binom{n}{k-j}$ zulke deelverzamelingen, en het
somprincipe over $j$ geeft de identiteit van Vandermonde.

Met $m = n = k$:
\[
\binom{2n}{n} = \sum_{j=0}^n \binom nj \binom{n}{n-j}
= \sum_{j=0}^n \binom nj^{2},
\]
waarbij de symmetrie $\binom{n}{n-j} = \binom nj$ gebruikt wordt.
\end{solution}

\begin{solution}{exo:g12:comb:10}
\emph{Via \cref{exo:g12:comb:7}:}
\[
\sum_{k=1}^n k\binom nk = \sum_{k=1}^n n\binom{n-1}{k-1}
= n\sum_{j=0}^{n-1}\binom{n-1}{j} = n\,2^{n-1},
\]
volgens \cref{cor:g12:comb:sums}. \emph{Via afleiden:} het afleiden van
$(1+x)^n = \sum_k \binom nk x^k$ geeft
$n(1+x)^{n-1} = \sum_k k \binom nk x^{k-1}$; vul $x = 1$ in.
\end{solution}

\begin{solution}{pb:g12:comb:1}
\textbf{1.} $26^2 \times 10^3 = 676\,000$ nummerplaten. BANANA: $6$ letters
met A verdrievoudigd en N verdubbeld: $\frac{6!}{3!\,2!} = 60$ anagrammen.

\textbf{2.} $\binom{32}{5} = 201\,376$ handen; $\binom42 \binom{28}{3} = 6
\times 3\,276 = 19\,656$ met precies twee azen.

\textbf{3.} Een pad is een woord met $4$ R's en $3$ B's: kies de plaatsen van
de B's: $\binom73 = 35$.

\textbf{4.} $(1+x)^4 = 1 + 4x + 6x^2 + 4x^3 + x^4$. Voor $x = 1$:
$\sum_k \binom nk = 2^n$; voor $x = -1$: $\sum_k (-1)^k \binom nk = 0$ ---
de rijsommen en de alternerende rijsommen van de \omterm{prop:g11:binom:pascal}{driehoek van Pascal}.

\textbf{5.} Comités van $k$ personen met een voorzitter, uit $n$: kies het
comité en daarna zijn voorzitter ($\binom nk \times k$), of de voorzitter en
daarna de overige leden ($n \times \binom{n-1}{k-1}$): gelijk. Sommeren over
$k$: het rechterlid sommeert tot $n \sum_j \binom{n-1}{j} = n\,2^{n-1}$.

\textbf{6.} Een \omterm{def:g12:seq:sequence}{rij} van $10$ sterren en $3$ strepen codeert de bestelling
(de bollen van smaak 1 vóór de eerste streep, enzovoort); de \omterm{def:g12:seq:sequence}{rij} telt $13$
symbolen en ligt vast door de plaatsen van de strepen:
$\binom{13}{3} = 286$ bestellingen.

\textbf{7.} $12$ sterren, $2$ strepen: $\binom{14}{2} = 91$.

\textbf{8.} Met $x', y', z' \geq 0$ en $x' + y' + z' = 9$:
$\binom{11}{2} = 55$.

\textbf{9.} Een monoom $a^i b^j c^k$ met $i + j + k = 5$: $\binom72 = 21$.

\textbf{10.} Met de formule: $3$ sterren, $1$ streep: $\binom41 = 4$; de
opsomming: $(3,0)$, $(2,1)$, $(1,2)$, $(0,3)$: overeenstemming.

\textbf{11.} ``Identiek'' kwam binnen op het ogenblik dat een bestelling niets
anders dan de \emph{aantallen} per smaak bleek te zijn --- de sterren dragen
geen namen. Worden de bollen in volgorde opgegeten, dan kiest elk van de $10$
verschillende plaatsen vrij een smaak: $4^{10} = 1\,048\,576$ mogelijkheden
--- een ander model en een andere wereld (\cref{met:g12:comb:model}: vraag
altijd \emph{geordend? verschillend? herhaling toegestaan?}).

\textbf{12.} $D_1 = 0$; $D_2 = 1$ (verwisselen); $D_3 = 2$ (de twee
$3$-cykels); $D_4 = 9$.

\textbf{13.} Gast 1 krijgt hoed $k \neq 1$: $n - 1$ keuzen. Krijgt gast $k$
hoed 1, dan verwisselen de overige $n - 2$ gasten hun eigen hoeden:
$D_{n-2}$ manieren. Krijgt gast $k$ hoed 1 \emph{niet}, herdoop dan hoed 1
tot de verboden hoed van gast $k$: de $n - 1$ resterende gasten verwisselen:
$D_{n-1}$ manieren. Dus $D_n = (n-1)(D_{n-1} + D_{n-2})$. Controle:
$D_4 = 3(2 + 1) = 9$; en $D_5 = 4(9 + 2) = 44$.

\textbf{14.} Trek van de $3! = 6$ toewijzingen die af welke minstens één hoed
op zijn plaats laten: er zijn er drie die een gegeven hoed vastleggen ($2!$
elk, $3 \times 2 = 6$), waarbij de paren dubbel geteld worden ($3$ paren,
$1!$ elk) en dus moeten terugkeren, en de identieke toewijzing ($1$) weer
afgetrokken wordt: $D_3 = 6 - 6 + 3 - 1 = 2$, dat wil zeggen
$3!\left(1 - 1 + \frac12 - \frac16\right) = 2$. In het algemeen is
$D_n = n!\sum_{k=0}^{n} \frac{(-1)^k}{k!}$.

\textbf{15.} $\frac{D_5}{120} = \frac{44}{120} \approx 0.3667$, al dicht bij
$\frac1\eu \approx 0.3679$: de alternerende som
$1 - 1 + \frac{1}{2!} - \frac{1}{3!} + \dots$ stapt naar $\eu^{-1}$. Op een
groot feest verwisselen de hoeden zich in ongeveer $36.8\,\%$ van de
gevallen --- de constante van de loterij en van de secretaresse, derde
waarneming.

\textbf{16.} $\P(\text{geldig}) = \frac{D_{10}}{10!} \approx 0.368$. Elke
hertrekking slaagt met \omterm{def:g10:proba:distribution}{kans} $\approx \frac1\eu$, dus het verwachte aantal
trekkingen is ongeveer $\eu \approx 2.7$: reken op drie rondes met de hoed.

\textbf{17.} Elke handdruk draagt $2$ bij aan de totale graadsom, dus de som
van alle handdrukaantallen van de gasten is even. Een som van \omterm{def:g10:numbers:sets}{gehele getallen}
is alleen even als het aantal oneven termen even is: de oneven schudders
komen in even aantallen. (Bij drie gasten: de mogelijke handdrukprofielen
hebben nooit precies één of drie oneven waarden --- ga de vier mogelijke
grafen na.)

\textbf{18.} $n = 1$: $1 = 1$. Geldt
$1^3 + \dots + n^3 = \left(\frac{n(n+1)}{2}\right)^2$, dan geeft $(n+1)^3$
erbij:
\[
\frac{n^2(n+1)^2}{4} + (n+1)^3
= \frac{(n+1)^2\left(n^2 + 4n + 4\right)}{4}
= \left(\frac{(n+1)(n+2)}{2}\right)^{\!2} :
\]
overerving. Voor $n = 3$: $1 + 8 + 27 = 36 = 6^2$.

\textbf{19.} Een pad naar $(n, n)$ zet $2n$ stappen en kruist de
antidiagonaal $x + y = n$ in precies één roosterpunt $(j, n - j)$; de eerste
helft is een pad met $j$ R's onder $n$ stappen ($\binom nj$ keuzen), en de
tweede helft, achterstevoren gelezen, eveneens ($\binom nj$ opnieuw, wegens
de symmetrie). Sommeren over het kruispunt geeft
$\binom{2n}{n} = \sum_j \binom nj^2$ --- de identiteit van Vandermonde,
getekend.

\textbf{20.} Vermenigvuldig fasen, tel gevallen op: de nummerplaten en de
pokerhanden. Codeer: paden als RB-woorden, bestellingen als sterren en
strepen. Tel twee keer: comités met voorzitter, handdrukken, doormidden
geknipte paden. Trek af en corrigeer: de verwisselde hoeden, met $\frac1\eu$
als restant. Volgende haltes: deze tellingen onder de breuken van de
kansrekening, en de padentellende machten van de verbindingsmatrices twee
hoofdstukken verderop.
\end{solution}
